\input{../header}

\begin{document}
%


%\onehalfspacing
\allowdisplaybreaks
%##################################################################
\section{PS\#11: Modeling with differential equations - \red{Answer key} }


\begin{enumerate}[leftmargin=0pt]

\item (\href{https://activecalculus.org/single/sec-7-5-modeling.html#ez-7-5-7}{AC 7.5 Exercise 7}) (Skydiver).
\begin{enumerate}
    \item Using $k$ to represent the constant of proportionality, write a differential equation that describes the rate of change of the skydiver's velocity.
    \begin{red}
        $\dfrac{dv}{dt} = 9.8 - k\cdot v.$
    \end{red}
    \item Find any equilibrium solutions and decide whether they are stable or unstable. Your result should depend on $k$.
    \begin{red} Stable eq. at $v = \dfrac{9.8}{k}$.
    \end{red}
    \item Suppose that the initial velocity is zero. Find the velocity $v(t)$.
    \begin{red}$v(t) = \dfrac{9.8}{k}\left(1 - \exp(-k\cdot t)\right)$
    \end{red}
    \item A typical terminal velocity for a skydiver falling face down is 54 meters per second. What is the value of $k$ for this skydiver?
    \begin{red} $k =\dfrac{9.8}{54} \approx 0.1815.$
    \end{red}
    \item How long does it take to reach 50\% of the terminal velocity?
    \begin{red} $t = \dfrac{\ln(2)}{0.1815} \approx 3.82\text{ seconds}$. \end{red}
\end{enumerate}

\item (\href{https://faculty.gvsu.edu/boelkinm/Home/ACS/sec-7-5-modeling.html#DfO}{AC 7.5 Exercise 8}) (Baby weight).
\begin{enumerate}
    \item Express this fact as a differential equation. \begin{red}
        $\dfrac{dW}{dt} = k\cdot\dfrac{1}{W}.$
    \end{red}
    \item Suppose that a baby weighs 8 pounds at birth and 9 pounds one month later. How much will he weigh at one year?
    \begin{red}
        $C = 64$, $k = 8.5$, $W(t) = \sqrt{17t+64}$, and $W(12) \approx 16.37$.
    \end{red}


\end{enumerate}


\item (\href{https://faculty.gvsu.edu/boelkinm/Home/ACS/sec-7-5-modeling.html#jmX}{AC 7.5 Exercise 9}) (Salt water tank).
\begin{enumerate}
    \item Since 3 gallons enter the tank every minute and 3 gallons leave every minute, what can you conclude about the volume of water in the tank?

    \begin{red}
        The volume of water in the tank is in fact constant.
    \end{red}
    \item How many grams of salt enter the tank every minute? \begin{red}
        $60$ grams / min
    \end{red}
    \item Suppose that $S(t)$ denotes the number of grams of salt in the tank in minute $t$. How many grams are there in each gallon in minute $t$?
    \begin{red}
        $\dfrac{S}{100}\ \dfrac{\text{grams}}{\text{gallon}}.$
    \end{red}
    \item Since water leaves the tank at 3 gallons per minute, how many grams of salt leave the tank each minute?
    \begin{red}
        $\dfrac{3S}{100}\ \dfrac{\text{grams}}{\text{minute}}.$
    \end{red}
    \item Write a differential equation that expresses the total rate of change of $S$. \begin{red}
        $\dfrac{dS}{dt} = 60 - \dfrac{3S}{100}.$
    \end{red}
    \item Identify any equilibrium solutions and determine whether they are stable or unstable.
    
    \begin{red}
        $S = 2000$.
    \end{red}
    \item Suppose that there is initially no salt in the tank. Find the amount of salt $S(t)$ in minute $t$.
    \begin{red}
        $S(t) = 2000 - 2000\cdot \exp\left(-\tfrac{3}{100}t\right)$
    \end{red}
\end{enumerate}

\end{enumerate}
\end{document}